\paragraph{Hint 1.} Use images of a finite generating set.

\paragraph{Hint 2.} Adjoin an inverse to $f$.

\paragraph{Hint 3.} Combine generating sets.

\paragraph{Hint 4.} Do not try to generate $I$.

\paragraph{Hint 5.} Use an infinitely generated ideal in a polynomial ring.

\paragraph{Hint 6.} Work over an affine open of the target.

\paragraph{Hint 7.} Affine inverse images are quasi-compact.

\paragraph{Hint 8.} Either use the tensor formula or the general theorem.

\paragraph{Hint 9.} Localize generators of an ideal contraction.

\paragraph{Hint 10.} Ideals upstairs correspond to ideals containing $I$.

\paragraph{Hint 11.} Induct on $n$.

\paragraph{Hint 12.} Present it as a quotient of a polynomial ring.

\paragraph{Hint 13.} Reduce modulo earlier variables.

\paragraph{Hint 14.} Use the components as an open cover.

\paragraph{Hint 15.} Use a finite union of finite affine covers.

\paragraph{Hint 16.} A fiber is a base change.

\paragraph{Hint 17.} Work affine-locally and combine algebra generators.

\paragraph{Hint 18.} Add quasi-compactness.

\paragraph{Hint 19.} Tensor the finite set of generators.

\paragraph{Hint 20.} Use the previous exercise plus quasi-compactness.

\paragraph{Hint 21.} Restrict a Noetherian affine cover and refine by distinguished opens.

\paragraph{Hint 22.} Use the topological characterization.

\paragraph{Hint 23.} Fields are Noetherian.

\paragraph{Hint 24.} Use infinitely many copies of any nonempty finite-type scheme.

\paragraph{Hint 25.} Affine line over a field suffices.

\paragraph{Hint 26.} Compare affine line and a finite set of points.
